\section{模型的评价与总结}
\label{sec:evaluation}

\subsection{模型的优点}
三问共用同一个基准和一套递进的接收域。三个模型都以同图整图单Task单核Makespan为加速比基准，逐步扩大张量的可复用范围：问题一按Task统计接收域，问题二按核心统计，问题三在核心驻留之上加入Cache查询时序。三问的结构搬运与容量约束公式均可逐项对应到评估程序的输出字段，问题二的分解恒等式在500个结果中逐行成立。

问题三用同一计划分别在B、C配置下评价，把Cache收益拆成硬件效应和调度适配两段（式\eqref{eq:cache-factors}）；五核100图中54图因硬件缩短工期、7图因调度适配进一步改善，两类效应可独立量化。继承方案在已评价计划中逐图取最短，三问五核均无慢于单核基准的图，相邻核数之间也没有退化。

\subsection{模型的不足}
16张只有一个弱连通分量的图，五核平均加速比仅1.17至1.30（表\ref{tab:topology-speedup}），切图能力是当前方法最弱的环节。主矩阵每个组合只对少数候选做完整评价，候选排序直接决定结果；预算实验（表\ref{tab:budget-sensitivity}）显示，增加评价次数可以继续提高部分图的加速比，但轻量评分与实际Makespan的相关性尚未检验。

问题三的调度适配收益很小：五核只有7图因调整计划而改善，平均适配比1.0011；六图在参数扰动下重优化全部选回原计划。各场景$F/LB$的第90百分位超过4.4（表\ref{tab:eval-cost}），说明大多数图的搜索空间仍有改善余地；问题一最大的单组合候选评价累计约1198 s，且没有端到端计时。全部工期由评估程序在CPU上的事件模拟给出，针对该评价函数进行了有限候选搜索。

\subsection{算法复杂度与计算成本}
算法成本分为图索引、候选构造与评分、完整评价三个阶段。令$|V|$为非COPY操作数，$|E_0|$为原始边数，$|E|$为去重后的计算依赖数，$|\mathcal T|$为张量数，$P$为逐张量枚举的生产者--消费者配对数。生产消费索引的代价为$O(|V|+|\mathcal T|+|E_0|+P)$，堆拓扑序及后继排序为$O((|V|+|E|)\log |V|)$，弱连通分量遍历为$O(|V|+|E|)$；释放优先序每步对全部就绪操作排序后取前8项，最坏为$O(|V|^2\log |V|)$。若生成$M$个候选、单候选检查与评分成本为$C_{\rm cand}$，候选阶段为$O(MC_{\rm cand}+M\log M)$。完整评价另需展开Task、COPY和物理版本，其成本随展开后的事件数$E_{\rm ev}$和容量扫描增长。每个图--核数--场景的常规完整调用额度为$\max\{\min(q,2),m+1\}$，其中$m$为去重后必评计划数；若最优当前慢于上一核继承方案，补评回退计划，额度不受此限制。

表\ref{tab:eval-cost}统计最终方案的下界距离与单次完整评价耗时，表\ref{tab:algorithm-timing}统计有计时字段的963条候选。$F/LB$的第90百分位超过4.4，说明大多数图的搜索空间仍有改善余地。
\begin{table}[H]\centering\small
\caption{最终方案的下界距离与单次评价耗时}\label{tab:eval-cost}
\setlength{\tabcolsep}{4pt}
\begin{tabular}{lrrrrr}\toprule
场景 & 组合数 & \shortstack{$F/LB$\\中位数} & \shortstack{$F/LB$\\第90百分位} & \shortstack{评价耗时\\中位数/s} & \shortstack{评价耗时\\第90百分位/s}\\\midrule
A & 500 & 1.325 & 4.617 & 1.484 & 53.487\\
B & 500 & 1.349 & 4.486 & 0.755 & 6.613\\
C & 500 & 1.314 & 4.486 & 0.752 & 6.618\\\bottomrule
\end{tabular}\end{table}
\begin{table}[H]\centering\small
\caption{有完整计时字段的候选评价耗时（秒）}
\label{tab:algorithm-timing}
\setlength{\tabcolsep}{2.4pt}
\begin{tabular}{lrrrrr}\toprule
场景 & 成功候选数 & \shortstack{展开检查\\中位数} & \shortstack{展开检查\\P95} & \shortstack{完整评价\\中位数} & \shortstack{完整评价\\P95/最大值}\\\midrule
A & 302 & 1.196 & 5.310 & 3.011 & 93.317 / 384.720\\
B & 275 & 1.578 & 8.684 & 1.680 & 9.315 / 22.891\\
C & 280 & 1.534 & 8.753 & 1.764 & 8.776 / 15.102\\\bottomrule
\end{tabular}
\end{table}
典型图的一次完整评价只需1至3 s，A场景大图的尾部则长得多。把同一组合的候选耗时相加，A/B/C的组合级中位数分别为3.99/2.42/2.32 s，P95为225.47/51.51/39.90 s，最大值为1198.43/211.71/212.42 s。对超大图，应按操作数和展开事件数设定候选上限，优先保证基准和继承方案被评价，再按剩余配额扩展搜索。

\subsection{模型改进与推广}
评价预算和代理质量是首要改进方向。按图规模和$F/LB$距离自适应分配评价次数，并设置端到端时间上限（以5至10 min为参考），可使每个用例在合理时间内完成；检验轻量评分与实际Makespan的秩相关，能指出代理在哪类图上失效，从而有针对性地改进排序依据。

单分量大图的加速有限，根源在于拓扑序切块对紧密耦合的长链并行化能力有限。对这类图，可按关键路径分段，在每段内同时平衡M、V两条流水线，并让相邻分段优先安排在同一核心，减少跨核等待。

模型适用于依赖图、操作成本和片上容量已知，且能对完整计划做事件评价的调度任务。结构搬运公式刻画接收域，物理副本的生存区间对应容量换出，FIFO查询与完成事件决定共享读服务；这一框架可以推广到其他多核加速器的图划分与调度问题，具体性能取决于目标硬件的事件规则。
